\section{预测与预警扩展：基于贸易网络的 SIRI 风险预测}

\subsection{扩展目的}

前文分析主要回答两个问题：一是中国半导体相关产品进口来源结构如何变化，二是这些结构变化如何反映在 SIRI 供应链风险指数及其解释性回归中。作为进一步扩展，本节将问题推进到预测层面：如果已经观察到某一年产品的进口来源网络和风险特征，能否对下一年的 SIRI 风险水平形成可检验的预测。这一部分不替代前文的描述统计和回归主线，而是检验 SIRI 是否可以从事后评价指标扩展为下一年风险预测任务。

图 \ref{fig:v03-workflow} 概括了从 BACI 贸易数据到下一年 SIRI 预测评估的流程。

\begin{figure}[H]
    \centering
    \includegraphics[width=0.98\textwidth]{v03_prediction_workflow.png}
    \caption{贸易网络预测扩展流程}
    \label{fig:v03-workflow}
\end{figure}

\subsection{数据扩展与图样本构造}

在数据与网络构造上，本节将产品范围从前文重点讨论的核心 4 个 HS6 产品扩展到 20 个半导体相关 HS6 产品。对每一个产品和每一个年份，本文构造一张中国进口来源网络：节点包括中国以及该年向中国出口该产品的来源国，边表示来源国向中国的出口关系，边权定义为该来源国在中国该产品年度进口额中的份额。因此，一个产品--年份对应一个年度图样本，图结构刻画的是进口来源集中度、来源国组合以及不同来源国在中国进口中的相对权重。样本规模与时间切分见表 \ref{tab:v03-overview}。

\begin{table}[H]
  \centering
  \caption{v0.3 贸易网络预测扩展样本概况}
  \label{tab:v03-overview}
  \begin{tabular}{lr}
    \toprule
    项目 & 数值 \\
    \midrule
    扩展产品数 & 20 \\
    年度图样本数 & 320 \\
    训练样本 & 240 \\
    验证样本 & 40 \\
    测试样本 & 40 \\
    GDELT 状态 & 未启用（BACI-only） \\
    \bottomrule
  \end{tabular}
\end{table}


\subsection{预测任务与模型比较}

预测设计采用一年期前瞻设定，即使用年份 $t$ 的贸易网络和风险特征预测 $t+1$ 年的产品级 SIRI。为避免把预测结果解释为政策冲击的因果识别，本文仅将其作为供应链风险预警的探索性任务，并在时间维度上划分训练、验证和测试样本。模型比较包括三类：Naive 基线使用上一年 SIRI 直接预测下一年 SIRI；Ridge 模型使用表格化风险特征进行线性正则化预测；贸易网络原型模型则利用图结构传播后的网络特征进行预测，用于检验年度进口来源图是否能提供额外的结构信息。这里的贸易网络原型对应轻量级 gcn\_numpy 原型，不等同于完整 GDELT-GCN 模型。测试集模型比较结果见表 \ref{tab:v03-metrics}。

\begin{table}[H]
  \centering
  \caption{v0.3 测试集预测指标比较}
  \label{tab:v03-metrics}
  \begin{tabular}{llrrrr}
    \toprule
    模型 & 样本范围 & MAE & RMSE & Spearman & 样本数 \\
    \midrule
    Naive & 全部20个产品 & 3.978 & 7.899 & 0.812 & 40 \\
    Naive & 核心4产品 & 1.667 & 2.117 & 0.881 & 8 \\
    Ridge & 全部20个产品 & 4.331 & 7.577 & 0.777 & 40 \\
    Ridge & 核心4产品 & 2.070 & 2.463 & 0.881 & 8 \\
    贸易网络原型 & 全部20个产品 & 61.179 & 263.756 & 0.601 & 40 \\
    贸易网络原型 & 核心4产品 & 1.974 & 2.444 & 0.857 & 8 \\
    \bottomrule
  \end{tabular}
  \caption*{注：当前 v0.3 结果为 BACI-only 贸易网络预测扩展，未接入真实 GDELT 新闻事件数据；“贸易网络原型”对应程序输出中的 gcn\_numpy 行，不代表完整 GDELT-GCN 模型。}
\end{table}


图 \ref{fig:v03-metrics-bar} 将测试集 MAE 和 RMSE 结果转化为分组柱状图。全部 20 个产品范围内，Naive 和 Ridge 的误差明显低于贸易网络原型；在核心 4 产品范围内，三类模型误差更接近，但复杂原型仍没有形成稳定优势。

\begin{figure}[H]
    \centering
    \includegraphics[width=0.96\textwidth]{prediction_metrics_bar.png}
    \caption{预测模型 MAE/RMSE 对比柱状图}
    \label{fig:v03-metrics-bar}
\end{figure}

\subsection{结果解读与预警价值}

结果显示，样本构造与评估流程表明，20 个产品、320 个年度图样本可以被组织为可复现的下一年 SIRI 预测任务。但从测试集指标看，在当前年度样本规模下，简单基线仍然具有较强竞争力，复杂贸易网络模型尚未体现出相对 Naive 或 Ridge 的稳定优势；其中贸易网络原型在全部 20 个产品上的误差明显偏大，说明该未约束原型在跨产品预测时仍不稳定。因此，贸易网络原型应被理解为探索性扩展，而不是对前文回归结论的强化证据。尤其需要说明的是，当前结果尚未启用真实 GDELT 新闻事件数据；本轮 BACI-only 扩展的主要价值在于建立产品--年份图样本、预测标签、基线比较和评估指标的统一框架，为后续接入 GDELT 新闻事件压力变量提供数据结构和评估框架。
